% !TeX root = ../papers/nc-p3.tex
\section{整な反対八次曲面の三重点と九余因子の接錐}
\label{sec:cofactor-tangent-nets}
\begin{theorem}\label{thm:cofactor-tangent-nets}
$D=V(\det N)$ を整な反対八次曲面、$Z$ を一般核の像の閉包とする。
$N(p)=yy^{\mathsf t}$ となる任意の点 $p$ で三次の接錐が非零なら、次が成り立つ。
\begin{enumerate}
\item $Z$ が曲線なら接錐は重複を許した三直線である。特に七・八余因子の場合は常にこの形になる。
\item 余因子が九次元の場合、その唯一の核二次曲面 $Q$ の平面 $y^\perp$ への制限の階数を $r$ とする。
$r=3$ なら接錐は三直線か、節点または尖点のある整三次曲線である。
整な場合、接行列の階数はその唯一の特異点で1である。
$r=2$ なら直線と二次式の積、$r=1$ なら平方一次因子を持つ三直線である。
\item 一つでも滑らかな三次接錐があれば、余因子十個は全て独立で、$Z$ は曲面である。
\end{enumerate}
\end{theorem}
\begin{proof}[接行列と核の特殊化]
$p$ で単元対角成分を消去し、残る対称三行の一次項を $S(v)$ と置く。$v$ は入力の三接座標である。
$d$ の三次項は $\det S(v)$ の非零定数倍であり、$\operatorname{adj}N$ の二次項の非零となり得る部分は
$y^\perp$ 上の $\operatorname{adj}S(v)$ に限られる。整な $D$ で核を消す二次式全体が
余因子の定数関係である。これを $y^\perp$ に制限する写像は単射である。核に入る二次式は平面の式で割れ、
二平面の和に整な非退化核像 $Z$ が入ることになるからである。

$\det S=0$ に次数2以上の既約成分 $C$ があれば、それは単純成分で一般点は $S$ の階数2点である。
その方向での核は $Z\cap y^\perp$ に入る。実際 $p$ を吹き上げたチャートで
$d(p+tv)/t^3$ の $t=0$ での式は $\det S(v)$ であり、$C$ の一般点で方向微分の一つが非零である。
従ってその点を $D$ の階数3点から特殊化でき、Schur補行列を $t$ で割った極限の核が $\ker S(v)$ となる。
$Z$ が曲線ならこの平面切断は有限で、$C$ 上の一般核は固定ベクトル $u$ になる。
$S(v)u$ の一次各成分が $C$ 上で消えれば恒等零なので $\det S\ne0$ に反する。
これで第一項が従う。
\end{proof}

\begin{proof}[九余因子の平面二次関係]
$r=3$ では座標を二項二次式の評価基底に取り、関係を
$(\operatorname{adj}S)_{02}=(\operatorname{adj}S)_{11}$ とする。
逆Bezout同型で $S$ を $\bigwedge^2 S_3$ に移すと関係はPlücker二次式そのものである。
実際六座標を $p_{ij}$ とすれば
\[
 S=\begin{pmatrix}p_{01}&p_{02}&p_{03}\\p_{02}&p_{03}+p_{12}&p_{13}\\p_{03}&p_{13}&p_{23}\end{pmatrix},
 \quad (\operatorname{adj}S)_{11}-(\operatorname{adj}S)_{02}
 =p_{01}p_{23}-p_{02}p_{13}+p_{03}p_{12}.
\]
$S(v)$ の線型包が三次元未満なら行列式は高々二変数の三次式で既に分解する。
三次元ならGrassmann多様体 $\operatorname{Gr}(2,S_3)$ 内の平面となる。
対応する射影直線族は各二本が交わるので、全てが一点を通るか、一平面に入る。
実際交わる二本を取ると、その張る平面外にある第三の直線はその交点を通り、
この三本と交わる全ての直線もそこを通る。この二場合は $c\wedge S_3$ と
$\bigwedge^2 W$（$\dim W=3$）である。

第一の場合の行列式は固定三次式 $c$ と動く三次式との結果式なので三直線となる。
第二の場合は $W$ に共通根があれば行列式が恒等零である。共通根がなければ
$\mathbb P^1\to\mathbb P(W^*)$ は基点なし次数3で、像は非退化である。
従って整三次曲線への正規化射である。$W$ 内の二次元部分空間の二基底の結果式が
零であることは、対応する $\mathbb P(W^*)$ の点がこの像にあることに等しい。
従ってその暗黙三次式が接錐そのものである。像の算術種数が1で幾何種数0ゆえ
節点か尖点を一つ持つ。そこでは対応する二切断が長さ2の逆像因子（尖点では二重根）を共有し、
二つの三次式のBezout行列の階数は1である。

$r=2$ なら関係を $S_{02}S_{12}=S_{01}S_{22}$ に取る。
$S_{22}\ne0$ のとき一次式の素因子分解で $S_{02}$ または $S_{12}$ が $S_{22}$ の定数倍となる。
前者を $c S_{22}$ とすれば $S_{01}=c S_{12}$ であり、第一列から最終列の $c$ 倍を引いて
同じ行操作をすると一行と二行のブロック和となる。後者も同様である。
$S_{22}=0$ なら二つの非対角成分の一方も零で、行列式は一次式と一次式の平方の積になる。
$r=1$ なら下二行二列の主小行列式が恒等零で、その対称一次行列の像は固定直線に入る。
これは二次Veronese円錐が射影直線を含まないことによる。そのブロックを対角化して展開すると
行列式は一次式と別の一次式の平方の積である。階数0の制限は先の単射性で起きない。
最後の主張は七未満の余因子の排除と以上の全場合から従う。
\end{proof}

\path{computations/verify_cofactor_tangent_nets.py} はPlücker二次式、固定三次式による三直線、
および有理三次曲線の節点での接行列の階数を記号検算する。
